% Adapte du code TikZ de Heritiana ; voir source-exp-ln-tcd.tex.
\begin{tikzpicture}[line cap=round,line join=round,>=stealth,x=.85cm,y=.85cm,font=\footnotesize]
                  \draw[->,color=bmtGris] (-4.2,0) -- (5.24,0);
                  \foreach \x in {-4,-3,-2,-1,1,2,3,4,5}
                  \draw[shift={(\x,0)},color=bmtGris] (0pt,2pt) -- (0pt,-2pt) node[below] {\footnotesize $\x$};
                  \draw[->,color=bmtGris] (0,-4.2) -- (0,5.3);
                  \foreach \y in {-4,-3,-2,-1,1,2,3,4,5}
                  \draw[shift={(0,\y)},color=bmtGris] (2pt,0pt) -- (-2pt,0pt) node[left] {\footnotesize $\y$};
                  \draw[color=bmtGris] (0pt,-10pt) node[right] {\footnotesize $0$};

  % Deux branches parametrees par le meme t : points exactement symetriques.
  \begin{scope}
    \clip (-4,-4) rectangle (5,5);
    \draw[bmtVetiver,line width=1.2pt,smooth,samples=220,domain=-4:ln(5)]
      plot ({exp(\x)},\x);
    \draw[bmtIndigo,line width=1.2pt,smooth,samples=220,domain=-4:ln(5)]
      plot (\x,{exp(\x)});
    \draw[bmtGris,dashed] (-4,-4)--(5,5);
  \end{scope}
  \node[bmtIndigo,anchor=west,fill=bmtPapier,inner sep=2pt] at (1.7,4.3) {$y=e^x$};
  \node[bmtVetiver,anchor=south,fill=bmtPapier,inner sep=2pt] at (3.8,1.7) {$y=\ln x$};
  \node[bmtGris,anchor=north west] at (3.3,3.2) {$y=x$};
\end{tikzpicture}
